\documentclass[11pt,a4paper]{article}
\usepackage[utf8]{inputenc}
\usepackage[T1]{fontenc}
\usepackage{lmodern,amsmath,amssymb,textcomp}
\usepackage[margin=24mm]{geometry}
\usepackage{longtable,booktabs,array,enumitem}
\usepackage[hidelinks]{hyperref}
\setlength{\parindent}{0pt}
\setlength{\parskip}{6pt}
\emergencystretch=3em
\begin{document}
{\LARGE\bfseries Arbitrarily long prime-exponent gaps in complete-interval subset products\par}\medskip
{\small Provisional mathematical authors: Rachel Teo, Wenyi Wang, Hamdi Abderrahmene\par}\medskip
{\small Judging and evaluation record: the submissions discussed in this document have been judged and evaluated by Alejandro Zarzuelo Urdiales for OpenMath 2026. This dated review records the stated verification, classification and unresolved holds; it is a preliminary evaluation rather than a signed award decision.\par}

{\small Event provenance: OpenMath 2026 ran 27 September-2 October 2026 with worldwide online participation. Detailed organizer listings specify Harvard Innovation Labs for the 27 September in-person event; the OpenMath programme advertises a hybrid kickoff at MIT. Sundai identifies its Harvard/MIT community. Organizer sources: \url{https://rsihouse.ai/openmath} ; \url{https://luma.com/yzp9abvr} ; \url{https://www.sundai.club/events/boston/recursive-learning-hack-with-harvard-innovation-labs}\par}

{\small Rachel Teo  -  Sakana AI.\par}

{\small Wenyi Wang  -  Sakana AI.\par}

{\small Hamdi Abderrahmene  -  Sakana AI.\par}

{\small Affiliations follow the recorded author declarations or public institutional/profile sources. Preferred author names, author order and publication affiliations await author review. This editorial compilation does not establish anthology coauthorship.\par}

{\small Original-result working manuscript | OpenMath 2026 | 5 October 2026\par}\medskip
{\small Central source and adjudication archive: \url{https://github.com/alejandrozu/openmath-2026-judging}\par}

\subsection{Abstract}
The distinct-subset-product interpolation question associated with OEIS A060957 is refuted. More strongly, for every positive gap length and every prescribed lower bound on a prime, the construction realizes an exact fibre consisting of two blocks of four attainable exponents separated by the prescribed missing block.

\subsection{The original interpolation problem}
Let P(n) be the set of products of subsets of \{1,...,n\}. A factor may be chosen once or omitted: this is not the unrestricted multiplicative monoid generated by that interval. The question asks whether, when m and p\textasciicircum{}a m are attainable for a prime p\ensuremath{\leq}n, all the intermediate p-exponent levels are attainable. Its complete-interval condition matters. A hole in a hand-picked collection of factors would not answer the question, because the omitted interval factors could fill that hole.

The submitted result handles all factors of the interval simultaneously. Its basic endpoint supplies m, p and n with m and p\textasciicircum{}5m in P(n), but p\textasciicircum{}4m outside P(n). The stronger endpoint specifies the entire exponent fibre, rather than just proving that one level is missing. For every g\ensuremath{\geq}1 and lower bound B, there are n,p,d,h with p prime, B<p\ensuremath{\leq}n, d>0 and p not dividing d, for which p\textasciicircum{}e d is attainable precisely at e=h,h+1,h+2,h+3,h+g+4,h+g+5,h+g+6,h+g+7. The g omitted levels are h+4 through h+g+3.

\subsection{Encoding and the exposed-face construction}
The proof represents integer products by their prime-exponent vectors, separating the distinguished p coordinate from all other coordinates. A graph-based atom construction controls the possible exponent increments. In the generalized construction, a circulant system decomposes the complete graph on an odd number of vertices into oriented spanning two-factors. Every vertex sees every other vertex once through the shifts and their inverses. This identity is proved from the circulant model; it is not left as an existence hypothesis in the final theorem.

Integral feasibility conditions in the height-one slice classify the allowable atoms. Summing inequalities by colour gives a degree-two budget; summing by vertex gives the complementary bound. These estimates constrain integral solutions to the two desired exponent blocks. The argument is algebraic and combinatorial rather than a search through arbitrary integer factors.

To transport that local construction to the complete interval, the proof chooses prime coordinates and an isolating rational linear functional. The positive-weight factors form a forced baseline W. Any subset product with the same p-free exponent vector must contain these positive-weight factors and can differ only through zero-weight atoms in the prescribed subspace. Consequently unrelated factors from \{1,...,n\} cannot silently create extra attainable levels. The complete-face theorem proves both directions of this classification. Finally, dividing the baseline by its p-primary factor produces d not divisible by p and an absolute exponent offset h.

\subsection{Formal target and proof obligations}
The important proof boundary is the passage from a model fibre to products of actual distinct interval factors. The selected Main and Generic.Final endpoints close that boundary through the complete-face, product-support and normalization modules. This publication pass independently replayed all 45 exact frozen authored source modules and the selected audit under Lean 4.33.1 and Mathlib 0df444a360eaa60ab8c11dca51a86af692955474, finishing on 5 October 2026 at 03:10:13 UTC. All four requested prints passed: the explicit subset-product counterexample, falsity of interpolation, exact arbitrary-gap eight-level fibre, and the consecutive-above corollary each depend only on propext, Classical.choice and Quot.sound. No selected endpoint uses admitted, native-evaluation or unknown axioms. Archived dual-version checks remain distinct from this fresh 45-source replay. Receipt: Verification/sakana-a060957-fresh-build.json.

The ordinary-proof expansion now gives the complete atom classification, the choice of prime parameters and separating linear functional, and the bridge preventing the remaining factors of the complete interval from filling the gap. The formal proof supplies a precise supplementary specification alongside that human-readable argument. No minimum counterexample size is asserted, and the theorem does not provide one fixed prime that works for every gap length.

\subsection{Prior work and significance}
Prime-coordinate encoding, exposed faces and semigroup holes are established methods. The proposed new contribution is their use to refute the complete-interval interpolation statement and to prescribe gaps of arbitrary length with an exact eight-level fibre. The closest recorded sources did not supply that theorem. This is a bounded priority comparison, so the paper should avoid an unqualified first-discovery claim until the author and referee checks are complete.

This family is a strong standalone number-theory/combinatorics preprint candidate. The qualitative counterexample answers the original question, while the arbitrary-gap theorem supplies the more durable mathematical story. The two statements should appear in one paper and receive one contribution-family treatment.

{\small This short manuscript records the construction and selected endpoints. The companion Original results anthology contains the extended ordinary proof exposition and its explicitly stated premises; the preserved Lean source remains the complete formal proof record.\par}

\subsection{Verification and reproducibility}
Fresh execution: sakana-a 060957: PASS; 45/45 custom modules compiled successfully; receipt Verification/sakana-a060957-fresh-build.json. Compiler pins, source hashes and endpoint axiom outputs are in Verification. Historical receipts and finite checks do not replace fresh replay.

\begin{samepage}
\subsection{SubsetProductGap.Generic.arbitrary\_gap\_eight\_levels\_above}
\begin{quote}\small\ttfamily theorem arbitrary\_gap\_eight\_levels\_above (g : \ensuremath{\mathbb{N}}) (hg : 1 \ensuremath{\leq} g) (lower : \ensuremath{\mathbb{N}}) :\newline     \ensuremath{\exists} n p d h : \ensuremath{\mathbb{N}}, p.Prime \ensuremath{\wedge} lower < p \ensuremath{\wedge} p \ensuremath{\leq} n \ensuremath{\wedge} 0 < d \ensuremath{\wedge} \ensuremath{\neg} p \ensuremath{\mid} d \ensuremath{\wedge}\newline       \ensuremath{\forall} e : \ensuremath{\mathbb{N}}, p\textasciicircum{}e * d \ensuremath{\in} productsOfSubsets n \ensuremath{\leftrightarrow}\newline         e = h \ensuremath{\vee} e = h + 1 \ensuremath{\vee} e = h + 2 \ensuremath{\vee} e = h + 3 \ensuremath{\vee}\newline           e = h + g + 4 \ensuremath{\vee} e = h + g + 5 \ensuremath{\vee} e = h + g + 6 \ensuremath{\vee} e = h + g + 7\end{quote}
{\small Source: proofs/a060957/OpHack/SubsetProducts/Generic/Final.lean:90.\par}

{\small Source SHA-256: a2e26af97d551980ecca1eb62f5db1362f75ada87f28795f99fe526651007d53\par}

\end{samepage}
\begin{samepage}
\subsection{SubsetProductGap.Generic.arbitrary\_gap\_consecutive\_above}
\begin{quote}\small\ttfamily theorem arbitrary\_gap\_consecutive\_above (g : \ensuremath{\mathbb{N}}) (hg : 1 \ensuremath{\leq} g) (lower : \ensuremath{\mathbb{N}}) :\newline     \ensuremath{\exists} n p M : \ensuremath{\mathbb{N}}, p.Prime \ensuremath{\wedge} lower < p \ensuremath{\wedge} p \ensuremath{\leq} n \ensuremath{\wedge} 0 < M \ensuremath{\wedge}\newline       M \ensuremath{\in} productsOfSubsets n \ensuremath{\wedge} p\textasciicircum{}(g + 1) * M \ensuremath{\in} productsOfSubsets n \ensuremath{\wedge}\newline         \ensuremath{\forall} j : \ensuremath{\mathbb{N}}, 1 \ensuremath{\leq} j \ensuremath{\to} j \ensuremath{\leq} g \ensuremath{\to} p\textasciicircum{}j * M \ensuremath{\notin} productsOfSubsets n\end{quote}
{\small Source: proofs/a060957/OpHack/SubsetProducts/Generic/Final.lean:113.\par}

{\small Source SHA-256: a2e26af97d551980ecca1eb62f5db1362f75ada87f28795f99fe526651007d53\par}

{\small\textbf{Sources and attribution:} \href{https://oeis.org/A060957}{1. oeis.org / A060957}; \href{https://arxiv.org/abs/2501.02695}{2. arXiv source: 2501.02695}; \href{https://github.com/alejandrozu/openmath-2026-judging/blob/d0ed487f487fa7ec814ff705ab55249983fe8707/OpenMath-Judging/review/entries/E02-a060957.txt}{3. alejandrozu/openmath-2026-judging / E02-a 060957.txt}; \href{https://github.com/alejandrozu/openmath-2026-judging}{Central source and adjudication archive}.\par}

\end{samepage}
\end{document}
